\documentclass{article}
\usepackage[T1]{fontenc}
\usepackage{lmodern}
\usepackage{graphicx}
\usepackage{amsmath,amssymb}
\usepackage[a4paper,margin=2cm]{geometry}
\usepackage[absolute,overlay]{textpos}
\setlength{\TPHorizModule}{1mm}
\setlength{\TPVertModule}{1mm}
\usepackage[indonesian]{babel}
\usepackage{hyperref}
\setlength{\emergencystretch}{2em}
% unit: Halaman 50

\begin{document}

% === Halaman 50 (llm) ===

\textbf{Contoh 4:}

\begin{tabular}{ll}
Plainteks & : \texttt{cryptoisshortforcryptography} \\
Kunci & : \texttt{abcdefabcdefabcdefabcdefabcd} \\
Cipherteks & : \texttt{CSASXTITUKWSTGQUCWYQVRKWAQJB} \\
\end{tabular}

\begin{itemize}
    \item Pada contoh di atas, \texttt{crypto} tidak dienkripsi menjadi kriptogram yang sama. 
    \item Mengapa bisa demikian?
\end{itemize}

\end{document}
